\documentclass[11pt]{ctexart}
\usepackage[a4paper,margin=1in]{geometry}
\usepackage{amsmath,amssymb,booktabs}
\usepackage{enumitem}
\usepackage{hyperref}

\title{Soft Half-t 不等式约束模型说明}
\author{}
\date{2026-06-05}

\begin{document}
\maketitle

\section{目标}

我们希望从一个候选阶段内的 feature 数据中判断该阶段是否存在单边不等式约束。给定一个阶段内的标准化 feature 样本
\[
\mathcal{X}=\{x_1,\ldots,x_n\},
\]
不等式约束有两种方向：
\[
\text{lower bound: } x\ge b,
\qquad
\text{upper bound: } x\le b.
\]

模型的核心思想是：如果某个阶段确实满足单边不等式约束，那么样本应该集中在某个边界 $b$ 的合法一侧，并且越靠近边界越常见；离边界越远，概率逐渐下降。同时，少量落在非法一侧的样本不应该导致 likelihood 直接变成零，因此使用 soft boundary 而不是 hard truncation。

\section{方向和 signed slack}

用一个方向变量 $d\in\{\mathrm{lower},\mathrm{upper}\}$ 表示不等式方向。定义符号
\[
 a_d=
 \begin{cases}
 +1, & d=\mathrm{lower},\\
 -1, & d=\mathrm{upper}.
 \end{cases}
\]

对任意样本 $x_i$，定义 signed slack：
\[
 y_i = a_d(x_i-b).
\]

于是：
\begin{itemize}[leftmargin=2em]
  \item $y_i\ge 0$ 表示样本在合法一侧；
  \item $y_i=0$ 表示样本正好在边界 $b$；
  \item $y_i<0$ 表示样本违反不等式约束。
\end{itemize}

lower bound 时，$y_i=x_i-b$；upper bound 时，$y_i=b-x_i$。

\section{Soft Half-t Density}

令 $t_\nu(z)$ 表示自由度为 $\nu$ 的标准 Student-t 密度：
\[
 t_\nu(z)
 =
 \frac{\Gamma\left(\frac{\nu+1}{2}\right)}
 {\sqrt{\nu\pi}\,\Gamma\left(\frac{\nu}{2}\right)}
 \left(1+\frac{z^2}{\nu}\right)^{-\frac{\nu+1}{2}}.
\]

Soft half-t 不等式模型定义为
\[
 p_{\mathrm{ineq}}(x\mid b,\tau,\nu,\alpha,d)
 =
 \begin{cases}
 \dfrac{1}{Z}\dfrac{2}{\tau}t_\nu\!\left(\dfrac{y}{\tau}\right),
 & y\ge 0,\\[1.2em]
 \dfrac{1}{Z}\dfrac{2}{\tau}t_\nu(0)\exp\!\left(\dfrac{y}{\alpha}\right),
 & y<0,
 \end{cases}
\]
其中
\[
 y=a_d(x-b).
\]

归一化常数为
\[
 Z = 1 + \frac{2}{\tau}t_\nu(0)\alpha.
\]

这个 density 可以理解为两部分：
\begin{itemize}[leftmargin=2em]
  \item 合法侧 $y\ge 0$ 是 half-Student-t 分布；
  \item 非法侧 $y<0$ 是指数泄漏项，用来给 violation 一个连续但快速衰减的概率。
\end{itemize}

\section{参数含义}

模型参数包括
\[
(b,\tau,\nu,\alpha,d).
\]

各参数的含义如下：
\begin{itemize}[leftmargin=2em]
  \item $b$：不等式边界。lower bound 中是下界，upper bound 中是上界。
  \item $d$：不等式方向，决定合法区域在 $b$ 的哪一侧。
  \item $\tau$：合法侧 half-t 的 scale，控制合法侧数据离边界的典型距离。
  \item $\nu$：Student-t 自由度，控制尾部厚度。$\nu$ 越小尾部越厚，模型越能容忍远离边界的合法侧样本。
  \item $\alpha$：soft boundary 的泄漏尺度，控制非法侧 violation 的惩罚强度。
\end{itemize}

通常令
\[
\alpha = \rho\tau,
\]
其中 $\rho>0$ 是相对 softness。$\rho$ 越小，非法侧衰减越快，模型越接近 hard inequality；$\rho$ 越大，非法侧越宽松。

\section{模型形状}

Soft half-t density 的最大值点在边界 $b$：
\[
\arg\max_x p_{\mathrm{ineq}}(x)=b.
\]

在合法侧，density 从 $b$ 开始随 $|x-b|$ 增大而下降：
\[
 p(x)\propto
 \left(1+\frac{(a_d(x-b)/\tau)^2}{\nu}\right)^{-\frac{\nu+1}{2}},
 \qquad a_d(x-b)\ge 0.
\]

在非法侧，density 从边界向外指数下降：
\[
 p(x)\propto \exp\left(\frac{a_d(x-b)}{\alpha}\right),
 \qquad a_d(x-b)<0.
\]

因此，这个分布是单峰、单边、边界锚定的。它不是对称分布，也不会退化成普通 baseline 分布；它始终表达“边界附近最高、合法侧 heavy tail、非法侧 soft penalty”的单边结构。

\section{边界和 Scale 的估计}

对于 lower bound，边界通常由低分位数估计：
\[
 b_{\mathrm{lower}} = Q_{q_b}(\mathcal{X}),
 \qquad q_b\approx 0.05.
\]

对于 upper bound，边界通常由高分位数估计：
\[
 b_{\mathrm{upper}} = Q_{1-q_b}(\mathcal{X}),
 \qquad q_b\approx 0.05.
\]

给定方向和边界后，先计算合法 slack：
\[
 s_i = \max(a_d(x_i-b),0).
\]

scale $\tau$ 用 slack 的高分位数估计，而不是用 second moment：
\[
 \tau_0
 =
 \frac{Q_q(\{s_i\})}{T_\nu^{-1}\left(\frac{1+q}{2}\right)},
 \qquad q\in[0.9,0.95].
\]
其中 $T_\nu^{-1}$ 是标准 Student-t CDF 的 inverse。

这个公式来自 half-t 分布的 CDF：
\[
 P(S\le s)=2T_\nu(s/\tau)-1.
\]
所以
\[
 Q_q(S)=\tau T_\nu^{-1}\left(\frac{1+q}{2}\right).
\]

使用高分位数估计 scale 的原因是：当大量样本贴近边界时，低分位数或 second moment 可能让 $\tau$ 过小，使模型变成过尖的 spike；高分位数可以让 scale 更多反映合法侧的整体 slack 宽度。

为了避免几乎常数的 feature 段产生无限尖锐的 density，还可以设置观测噪声下限：
\[
 \tau = \max(\tau_0,\epsilon),
\]
其中 $\epsilon$ 表示最小可辨别 feature 噪声尺度。

\section{Likelihood}

给定一个候选阶段的数据 $\mathcal{X}$，soft half-t 不等式模型的平均 negative log-likelihood 为
\[
 \mathrm{NLL}_{\mathrm{ineq}}
 =
 -\frac{1}{n}\sum_{i=1}^n
 \log p_{\mathrm{ineq}}(x_i\mid b,\tau,\nu,\alpha,d).
\]

如果需要比较 lower 和 upper 两个方向，则分别计算
\[
 \mathrm{NLL}_{\mathrm{lower}},
 \qquad
 \mathrm{NLL}_{\mathrm{upper}},
\]
并选择 likelihood 更好的方向：
\[
 \hat d = \arg\min_{d\in\{\mathrm{lower},\mathrm{upper}\}}
 \mathrm{NLL}_{\mathrm{ineq}}(d).
\]

\section{Baseline Comparison}

为了判断一个阶段是否真的表现出不等式结构，需要和一个无约束 baseline 分布比较。baseline 的作用不是给分段提供奖励，而是提供一个参照：只有当单边不等式模型明显优于无约束模型时，才认为该阶段存在不等式约束。

baseline 使用普通 Student-t 分布：
\[
 p_{\mathrm{base}}(x\mid \mu_b,\sigma_b,\nu_b)
 =
 \frac{1}{\sigma_b}
 t_{\nu_b}\!\left(\frac{x-\mu_b}{\sigma_b}\right).
\]

其平均 negative log-likelihood 为
\[
 \mathrm{NLL}_{\mathrm{base}}
 =
 -\frac{1}{n}\sum_{i=1}^n
 \log p_{\mathrm{base}}(x_i\mid \mu_b,\sigma_b,\nu_b).
\]

Student-t baseline 可以用稳健估计或 MLE 估计。概念上，baseline 应该描述“这个阶段的 feature 只是一个普通单峰无约束分布”的解释能力。它不包含边界 $b$，也不区分 lower/upper 方向。

如果使用 MLE baseline，可以固定自由度 $\nu_b$，优化
\[
 (\hat\mu_b,\hat\sigma_b)
 =
 \arg\min_{\mu,\sigma}
 \left[-\frac{1}{n}\sum_{i=1}^n
 \log p_{\mathrm{base}}(x_i\mid \mu,\sigma,\nu_b)\right].
\]

为了避免某些低方差阶段产生异常大的 likelihood，baseline scale 和 inequality scale 可以共享同一个最小噪声尺度：
\[
 \sigma_b \leftarrow \max(\sigma_b,\epsilon),
 \qquad
 \tau \leftarrow \max(\tau,
 \epsilon).
\]

这样做的含义是：feature 的测量或标准化后数值存在一个最小可分辨尺度，任何模型都不应该把 scale 缩到低于这个尺度。

\section{不等式方向选择}

对于每个候选阶段和每个 feature，可以分别拟合 lower-bound 和 upper-bound 两个 soft half-t 模型。

lower-bound 模型：
\[
 d=\mathrm{lower},
 \qquad
 b_{\mathrm{lower}}=Q_{q_b}(\mathcal{X}),
 \qquad
 y_i=x_i-b_{\mathrm{lower}}.
\]

upper-bound 模型：
\[
 d=\mathrm{upper},
 \qquad
 b_{\mathrm{upper}}=Q_{1-q_b}(\mathcal{X}),
 \qquad
 y_i=b_{\mathrm{upper}}-x_i.
\]

分别计算两个方向的不等式 NLL：
\[
 \mathrm{NLL}_{\mathrm{lower}}
 =
 -\frac{1}{n}\sum_i
 \log p_{\mathrm{ineq}}(x_i\mid b_{\mathrm{lower}},\tau_{\mathrm{lower}},\nu,\alpha_{\mathrm{lower}},\mathrm{lower}),
\]
\[
 \mathrm{NLL}_{\mathrm{upper}}
 =
 -\frac{1}{n}\sum_i
 \log p_{\mathrm{ineq}}(x_i\mid b_{\mathrm{upper}},\tau_{\mathrm{upper}},\nu,\alpha_{\mathrm{upper}},\mathrm{upper}).
\]

然后用相同的 baseline 作为参照，得到两个方向的 gain：
\[
 \mathrm{gain}_{\mathrm{lower}}
 =
 \mathrm{NLL}_{\mathrm{base}}-\mathrm{NLL}_{\mathrm{lower}},
\]
\[
 \mathrm{gain}_{\mathrm{upper}}
 =
 \mathrm{NLL}_{\mathrm{base}}-\mathrm{NLL}_{\mathrm{upper}}.
\]

如果方向未知，则选择 gain 最大的方向：
\[
 \hat d
 =
 \arg\max_{d\in\{\mathrm{lower},\mathrm{upper}\}}
 \mathrm{gain}_d.
\]

等价地，如果定义
\[
 \mathrm{score}_d
 =
 \mathrm{NLL}_{\mathrm{ineq},d}-\mathrm{NLL}_{\mathrm{base}},
\]
则选择 score 最小的方向：
\[
 \hat d
 =
 \arg\min_{d\in\{\mathrm{lower},\mathrm{upper}\}}
 \mathrm{score}_d.
\]

如果先验已经指定方向，例如某个 feature 只能是 upper-bound 或只能是 lower-bound，则不做方向搜索，只评估指定方向。

\section{激活判断}

方向选定后，判断该候选阶段是否激活不等式约束。设选定方向为 $\hat d$，对应的不等式 NLL 为
\[
 \mathrm{NLL}_{\mathrm{ineq}}
 =
 \mathrm{NLL}_{\mathrm{ineq},\hat d}.
\]

定义 likelihood gain：
\[
 \mathrm{gain}
 =
 \mathrm{NLL}_{\mathrm{base}}-\mathrm{NLL}_{\mathrm{ineq}}.
\]

或者定义 score：
\[
 \mathrm{score}
 =
 \mathrm{NLL}_{\mathrm{ineq}}-\mathrm{NLL}_{\mathrm{base}}
 = -\mathrm{gain}.
\]

激活规则为
\[
 \mathrm{active}
 =
 \mathbf{1}\left[\mathrm{gain}>\gamma\right],
\]
其中 $\gamma$ 是激活阈值。等价地，若使用 score，则
\[
 \mathrm{active}
 =
 \mathbf{1}\left[\mathrm{score}< -\gamma\right].
\]

这个判断的意义是：不等式模型必须比普通 Student-t baseline 好出一个 margin，才认为该阶段真的存在单边约束。否则，即使 soft half-t 本身能拟合数据，也不把它当作有效约束。

\section{分段中的 Likelihood / Cost}

在分段时，需要决定一个候选阶段是否因为满足不等式而获得约束项。这里区分两种情况。

如果不等式没有激活：
\[
 \mathrm{active}=0,
 \qquad
 C_{\mathrm{ineq}}=0.
\]

也就是说，未激活时不使用 baseline likelihood 影响分段。这样可以避免算法退化成“选择 baseline likelihood 最大的切分点”。baseline 只用于判断是否激活，不作为 inactive segment 的分段目标。

如果不等式激活：
\[
 \mathrm{active}=1,
\]
则使用 soft half-t 不等式模型的 total negative log-likelihood 作为该阶段的不等式代价：
\[
 C_{\mathrm{ineq}}
 =
 \lambda_{\mathrm{ineq}}\sum_{i=1}^{n}
 \left[-\log p_{\mathrm{ineq}}(x_i\mid b,\tau,\nu,\alpha,\hat d)\right].
\]

等价地，使用平均 NLL 表示为
\[
 C_{\mathrm{ineq}}
 =
 \lambda_{\mathrm{ineq}}\,n\,\mathrm{NLL}_{\mathrm{ineq}}.
\]

这里 $n$ 是候选阶段长度，$\lambda_{\mathrm{ineq}}$ 是不等式约束项权重。使用 total NLL 的原因是：分段目标应该考虑该阶段所有样本对约束的支持程度，而不是只看平均值。长阶段中稳定满足约束比短阶段中的偶然贴边更有统计意义。

如果整个分段方案包含多个阶段和多个 feature，则总不等式代价可以写成
\[
 C_{\mathrm{ineq,total}}
 =
 \sum_{k}\sum_{m}
 \mathbf{1}[\mathrm{active}_{k,m}=1]\,
 \lambda_{\mathrm{ineq}}
 \sum_{i\in \mathcal{S}_k}
 \left[-\log p_{\mathrm{ineq}}(x_{i,m}\mid b_{k,m},\tau_{k,m},\nu,\alpha_{k,m},d_{k,m})\right],
\]
其中 $\mathcal{S}_k$ 表示第 $k$ 个阶段包含的时间点集合，$m$ 表示 feature index。

\section{整体流程总结}

对每个候选阶段和每个 feature，方法流程如下：
\begin{enumerate}[leftmargin=2em]
  \item 取该阶段内的标准化 feature 样本 $\mathcal{X}$。
  \item 拟合无约束 Student-t baseline，得到 $\mathrm{NLL}_{\mathrm{base}}$。
  \item 拟合 lower soft half-t，得到 $\mathrm{NLL}_{\mathrm{lower}}$ 和 $\mathrm{gain}_{\mathrm{lower}}$。
  \item 拟合 upper soft half-t，得到 $\mathrm{NLL}_{\mathrm{upper}}$ 和 $\mathrm{gain}_{\mathrm{upper}}$。
  \item 如果方向未知，选择 gain 最大的方向；如果方向已知，只使用指定方向。
  \item 若最佳 gain 超过激活阈值 $\gamma$，则该 feature 在该阶段激活不等式约束。
  \item 激活时，分段目标中加入 soft half-t total NLL；未激活时，该 feature 的不等式 cost 为 0。
\end{enumerate}

\section{直观解释}

Soft half-t 不等式模型表达的是以下假设：

\begin{quote}
如果某个阶段存在单边约束，那么该阶段的 feature 应该在某个边界附近形成单峰密度；边界合法侧允许 heavy-tailed spread，非法侧允许少量 violation，但 violation 的概率随越界距离指数衰减。
\end{quote}

这个模型相比 hard truncated distribution 更稳定，因为 violation 不会导致 likelihood 变成 $-\infty$。相比普通 Student-t baseline，它又保留了明确的单边 inductive bias：最高密度固定在边界上，而不是自由地把中心放在样本均值或中位数附近。

\section{模型的主要优点和局限}

优点：
\begin{itemize}[leftmargin=2em]
  \item 明确表达 lower-bound 或 upper-bound 的单边结构。
  \item mode 锚定在边界 $b$，适合描述“贴边”的不等式约束。
  \item 合法侧使用 Student-t heavy tail，对合法但远离边界的点较稳健。
  \item 非法侧使用 soft exponential leakage，可以容忍少量 violation。
  \item 通过和 baseline 比较，可以避免把所有阶段都强行解释成不等式约束。
\end{itemize}

局限：
\begin{itemize}[leftmargin=2em]
  \item 如果一个多模态分布在某个方向上也能被解释成“边界附近高密度 + 合法侧 spread”，模型仍可能给出较高 likelihood。
  \item scale 太小时 likelihood 会异常尖锐，因此需要合理的 scale estimation 或 noise floor。
  \item 激活阈值控制模型的保守程度；阈值太低会漏检，太高会误检。
  \item 模型本身只描述单个 feature 的边缘分布，不直接建模多个 feature 之间的联合依赖。
\end{itemize}

\end{document}
